\appendix

\section{Dataset Details}
\label{sec:app_dataset_details}

We will provide more data statistics, detailed analysis, and sample display of \TVD.

\subsection{Genre Tag System}
\label{sec:app_genre_tag_system}

\begin{table*}[h]
    \centering
    \scriptsize
    \begin{tabularx}{\linewidth}{l|X}
    \toprule
    \makecell[c]{1st layer tag} & \makecell[c]{2nd layer tags} \\
    \midrule
    \begin{CJK*}{UTF8}{gbsn}奇幻\end{CJK*}/fantasy & \begin{CJK*}{UTF8}{gbsn}奇幻魔法\end{CJK*}/fantasy magic, \begin{CJK*}{UTF8}{gbsn}奇幻爱情\end{CJK*}/fantasy romance, \begin{CJK*}{UTF8}{gbsn}都市奇幻\end{CJK*}/urban fantasy \\
    \midrule
    \begin{CJK*}{UTF8}{gbsn}都市\end{CJK*}/urban & \begin{CJK*}{UTF8}{gbsn}都市奇幻\end{CJK*}/urban fantasy, \begin{CJK*}{UTF8}{gbsn}都市喜剧\end{CJK*}/urban comedy \\
    \midrule
    \begin{CJK*}{UTF8}{gbsn}喜剧\end{CJK*}/comedy & \begin{CJK*}{UTF8}{gbsn}都市喜剧\end{CJK*}/urban comedy, \begin{CJK*}{UTF8}{gbsn}恶搞闹剧\end{CJK*}/spoof farce, \begin{CJK*}{UTF8}{gbsn}黑色喜剧\end{CJK*}/dark comedy, \begin{CJK*}{UTF8}{gbsn}爱情喜剧\end{CJK*}/romantic comedy, \begin{CJK*}{UTF8}{gbsn}情景喜剧\end{CJK*}/situation comedy, \begin{CJK*}{UTF8}{gbsn}家庭喜剧\end{CJK*}/family comedy, \begin{CJK*}{UTF8}{gbsn}生活喜剧\end{CJK*}/life comedy, \begin{CJK*}{UTF8}{gbsn}搞笑\end{CJK*}/funny, \begin{CJK*}{UTF8}{gbsn}趣味\end{CJK*}/spice \\
    \midrule
    \begin{CJK*}{UTF8}{gbsn}惊悚\end{CJK*}/thriller & \begin{CJK*}{UTF8}{gbsn}悬疑\end{CJK*}/suspense, \begin{CJK*}{UTF8}{gbsn}惊悚\end{CJK*}/thriller \\
    \midrule
    \begin{CJK*}{UTF8}{gbsn}犯罪\end{CJK*}/crime & \begin{CJK*}{UTF8}{gbsn}高智商犯罪\end{CJK*}/high-IQ crime, \begin{CJK*}{UTF8}{gbsn}警匪犯罪\end{CJK*}/police and bandits crime, \begin{CJK*}{UTF8}{gbsn}烧脑罪案\end{CJK*}/brain burning crime, \begin{CJK*}{UTF8}{gbsn}侦探片\end{CJK*}/detective \\
    \midrule
    \begin{CJK*}{UTF8}{gbsn}科幻\end{CJK*}/sci-fi & \begin{CJK*}{UTF8}{gbsn}科幻动作\end{CJK*}/sci-fi action, \begin{CJK*}{UTF8}{gbsn}软科幻\end{CJK*}/soft sci-fi, \begin{CJK*}{UTF8}{gbsn}超自然\end{CJK*}/supernatural \\
    \midrule
    \begin{CJK*}{UTF8}{gbsn}动作\end{CJK*}/action & \begin{CJK*}{UTF8}{gbsn}科幻动作\end{CJK*}/sci-fi action, \\
    \midrule
    \begin{CJK*}{UTF8}{gbsn}爱情\end{CJK*}/romance & \begin{CJK*}{UTF8}{gbsn}奇幻爱情\end{CJK*}/fantasy romance, \begin{CJK*}{UTF8}{gbsn}熟龄爱情\end{CJK*}/mature romance, \begin{CJK*}{UTF8}{gbsn}曲折爱情\end{CJK*}/tortuous romance, \begin{CJK*}{UTF8}{gbsn}浪漫爱情\end{CJK*}/romantic love, \begin{CJK*}{UTF8}{gbsn}多角恋\end{CJK*}/polygonal love \\
    \midrule
    \begin{CJK*}{UTF8}{gbsn}题材\end{CJK*}/theme & \begin{CJK*}{UTF8}{gbsn}女性题材\end{CJK*}/female subject, \begin{CJK*}{UTF8}{gbsn}金融题材\end{CJK*}/finance, \begin{CJK*}{UTF8}{gbsn}律政题材\end{CJK*}/legal, \begin{CJK*}{UTF8}{gbsn}剧情故事\end{CJK*}/plot story, \begin{CJK*}{UTF8}{gbsn}行业剧\end{CJK*}/family comedy, \begin{CJK*}{UTF8}{gbsn}职场剧\end{CJK*}/workplace, \begin{CJK*}{UTF8}{gbsn}广告片\end{CJK*}/advertising \\
    \midrule
    \begin{CJK*}{UTF8}{gbsn}青春\end{CJK*}/teen & \begin{CJK*}{UTF8}{gbsn}青春\end{CJK*}/teen, \begin{CJK*}{UTF8}{gbsn}校园\end{CJK*}/school \\
    \midrule
    \begin{CJK*}{UTF8}{gbsn}情感\end{CJK*}/emotion & \begin{CJK*}{UTF8}{gbsn}父爱\end{CJK*}/paternal love, \begin{CJK*}{UTF8}{gbsn}婚姻生活\end{CJK*}/matrimony, \begin{CJK*}{UTF8}{gbsn}亲情\end{CJK*}/family affection \\
    \midrule
    \begin{CJK*}{UTF8}{gbsn}美剧\end{CJK*}/american & \begin{CJK*}{UTF8}{gbsn}美剧\end{CJK*}/american, \begin{CJK*}{UTF8}{gbsn}超级英雄\end{CJK*}/superhero \\
    \midrule
    \begin{CJK*}{UTF8}{gbsn}恐怖\end{CJK*}/horror & \begin{CJK*}{UTF8}{gbsn}恐怖\end{CJK*}/horror \\
    \midrule
    \begin{CJK*}{UTF8}{gbsn}冒险\end{CJK*}/adventure & \begin{CJK*}{UTF8}{gbsn}冒险\end{CJK*}/adventure \\
    \midrule
    \begin{CJK*}{UTF8}{gbsn}纪录片\end{CJK*}/documentary & \begin{CJK*}{UTF8}{gbsn}伪纪录片\end{CJK*}/mockumentary, \begin{CJK*}{UTF8}{gbsn}野史传奇\end{CJK*}/wild history, \begin{CJK*}{UTF8}{gbsn}人物传记\end{CJK*}/biography, \begin{CJK*}{UTF8}{gbsn}人物故事\end{CJK*}/character story \\
    \midrule
    \begin{CJK*}{UTF8}{gbsn}其他\end{CJK*}/other & \begin{CJK*}{UTF8}{gbsn}单元剧\end{CJK*}/unit play, \begin{CJK*}{UTF8}{gbsn}漫画改编\end{CJK*}/comic adaptation \\
    \bottomrule
    \end{tabularx}
    \caption{\label{tab:genre_tags}
    The hierarchical genre tags in \TVD.
    }
\end{table*}

Tab.~\ref{tab:genre_tags} lists the two layers of genre tags respectively in \TVD.
It can be found that viewers of TV drama source websites pay more attention to TV dramas with genre tags such as romance, crime, and fantasy.
At the same time, due to viewer preference, copyright purchase, and some hidden factors, some genre tags (e.g., documentary and adventure) are less, resulting in the imbalance of genre tags in \TVD.
Furthermore, when jointly considering the statistical results by genre tags of plots and TV dramas, it can be found that the proportion of plots in sci-fi and urban TV dramas is large, while the proportion of plots in comedy and theme TV dramas is small.
It shows that the plot density of TV dramas with different genre tags exists a gap, which is usually determined by the shooting methods and prominent objectives.

\subsection{Plot Clip}
\label{sec:app_plot_clip}

\begin{figure}[ht]
\centering
\includegraphics[width=1.0\textwidth]{figure/fig_duration.pdf}
\caption{
The distribution of duration of plot clip in \TVD.
}
\label{fig:duration}
\end{figure}

In \TVD, the core of parallel multi-modal data is the plot.
In Tab.~\ref{fig:duration}, we have counted the video clip duration in \TVD.
It shows the duration of the plot clip is mainly distributed between 0-2 minutes, which is significantly longer than existing datasets.
In addition, the plot clip can not only cover more long-distance dependent information, but also ensure the integrity of the plot.


\subsection{Bullet Screen Comments}
\label{sec:app_bsc}

\begin{figure}[ht]
\centering
\includegraphics[width=1.0\textwidth]{figure/fig_bsc_example.pdf}
\caption{
An example of BSC of the online video website.
}
\label{fig:bsc_example}
\end{figure}

As shown in Fig.~\ref{fig:bsc_example}, the BSC, as a real-time comment generated by viewers when watching a movie and TV dramas, will comment on the video content from various perspectives.
In particular, viewers will form a single round of opinion discussion with other viewers across time and space while watching TV dramas.
In addition, viewers express their agreement with the target BSC by clicking the `likes' button, which is also a simplified exchange of viewer opinions.
Therefore, the BSC data often contain a large amount of valuable behavior data, which can be used in commercial applications such as hot spot detection, topic mining, and movie review generation.

\subsubsection{Data statistics}
\begin{figure}[ht]
\centering
\includegraphics[width=0.7\textwidth]{figure/fig_bs_length.pdf}
\caption{
The length of the BSCs in \TVD.
}
\label{fig:bs_length}
\end{figure}

\begin{figure}[t]
\centering
\includegraphics[width=0.7\textwidth]{figure/fig_like.pdf}
\caption{
The distribution of user `likes' of the BSC in \TVD.
}
\label{fig:likes}
\end{figure}

In Fig.~\ref{fig:bs_length} and Fig.~\ref{fig:likes}, we have made detailed statistics on the information of the BSCs, including the length of the BSCs, the number of 'likes', and so on.
Among them, by analyzing the distribution of the length of the BSC, we can find that as an instant opinion, the BSC often does not carry too much content, and users often express their opinions in a short and concise way.
Besides, the statistical results of the number of 'likes' show that users have strong opinions to interact with other users when browsing, and the proportion of highly recognized BSCs is not low, which proves the value of BSC text laterally.

\subsubsection{Bullet screen comment and plot}
\label{sec:bsc_plot}
As real-time comments are sent by viewers when they watch videos, the number of BSCs can usually reflect the user's attention to specific sections in the video. 
Furthermore, these segments are often the key plots in TV dramas.
Therefore, we try to analyze the relationship between the popularity of BSCs in the video segment and the plot clip annotation information in \TVD.
We first use the sliding window to count the number of BSCs in intervals of 10 seconds, and then use the peak-detecting algorithm~\citep{peakdet2006bio} to obtain the top 5 hot segments (i.e., key plots) in each episode of TV dramas after smoothing the data, as shown in Fig.~\ref{fig:bsc_example}.
After obtaining these key plot clips, we use the plot clip annotation information in \TVD for matching.
When the overlapping rate of the annotated plot interval and the hotspot interval exceeds 80\%, it can be considered as hitting the key plot, the coverage rate of the key plot in \TVD is 83.3\%; when the overlap rate is set to 50\%, the key plot in \TVD The coverage rate is 97.7\%.

\begin{table*}
\centering
\small
\begin{tabular}{lccccccc}
\toprule
 & MLE & Miller\_Madow & Jackknife & Horvitz\_Thompson & Chao\_Shen & Wolpert\_Wolf & NSB \\
\midrule
plot & 7.05 & 7.07 & 7.07 & 7.12 & 7.09 & 7.18 & 7.09 \\
BSC & 6.64 & 6.65 & 6.65 & 6.71 & 6.68 & 6.73 & 6.66 \\
\bottomrule
\end{tabular}
\caption{\label{tab:entropy}
The entropy of plot and BSC text in \TVD. 
The units of values in the table are nat, 1 nat = $1/\mathrm{ln}2 \approx$ 1.44 bits.
}
\end{table*}

As the text was obtained by manual editing, the plot text and the BSC text are very different in expression habits, forms, and attention perspectives.
This is mainly due to the fact that plot annotators have stronger professionalism and focus on necessary narrative plot information.
However, the BSC comes from the viewers. 
Although the language usage of BSC is random, the pattern is uniform, and the focus points are also relatively consistent, focusing on actors and specific plots.
In order to analyze and then measure this difference, we wish to fully study the difference in the distribution of languages in these two texts (i.e., plot text and BSC text).
Specifically, referring to the work of~\citep{entropy2022acl}, we used various entropy estimators to calculate the plot and BSC text respectively (20,000 pieces each), and the results are shown in Tab.~\ref{tab:entropy}.

Specifically, each text entropy is calculated as follows:
\begin{itemize}
    \item Maximum Likelihood Estimation (MLE): This is the usual naive way to calculate the entropy of a sample--just measure the observed probabilities over outcomes and plug those in as if they were the true values (i.e. we have a model of the distribution that maximises the likelihood of our observations).
    \item Miller-Madow: Applies a first-order Taylor correction to MLE.
    \item Jackknife: Using subsamples of our sample, it attempts to project the entropy estimate for the whole sample.
    \item Horvitz-Thompson: Takes into account the likelihood of an outcome occurring at all in our sample--less-frequent outcomes contribute more to the entropy in order to compensate for MLE's negative bias.
    \item Chao-Shen: Modifies Horvitz-Thompson to use a coverage-reduced estimate of the probability of the outcomes. Empirically better than it.
    \item Wolpert-Wolf: First Bayesian entropy estimator with a Dirichlet prior.
    \item Nemenman-Shafee-Bialek (NSB): Improves upon Wolpert-Wolf.
\end{itemize}

From the results shown in Tab.~\ref{tab:entropy}, we can observe that the entropy of BSC text is consistently much lower than that of episodic text across all estimators.
In particular, when we select much more BSC texts (2,000,000 pieces) to measure entropy, the obtained results are basically consistent with the existing results (20,000 pieces).
We believe that this intuitively indicates that the internal language usage and expression of the BSC text tend to be more unified than the plot text, reflecting that users often have more fixed words or formats when posting BSCs.
Compared with BSCs, plot texts need to pay more attention to narrative quality, and the ways of words and sentences will be more diverse.

\section{Experiment Setup}
\label{sec:app_experiment_setup}

We will describe the dataset splitting, parameters, and two-stage (pre-training and fine-tuning) process of the framework in this section.


\subsection{Dataset splitting}
\label{sec:app_dataset_splitting}

We split training, validation, and test data in 8:1:1 for all downstream tasks (i.e., genre classification, cross-modal retrieval, and plot text ).
In addition, to ensure that entities, such as characters, appearing in the validation and test sets are already present in the training set, we make the plot clips of each TV drama appear uniformly but randomly in the training, validation, and test sets.


\subsection{Pre-training settings}
\label{sec:app_pretrain_setting}

In the three experiments, we used vision and text encoders are vit\_base\_p16\_224\footnote{\url{https://huggingface.co/google/vit-base-patch16-224}} and bert-base-chinese\footnote{\url{https://huggingface.co/bert-base-chinese}} respectively.
In the image-text pre-training model, we use the AdamW optimizer~\citep{adamw2017axriv} with the learning rate set to 1e-5 for the pre-training model.
The warm-up ratio is set to 10\%, and the learning rate decays linearly to 0 after 10\% of the total training steps.
The batch size, hidden size, and number of heads are set to 4096, 768, and 12, respectively.
We pre-trained our model for 100K steps on 8 NVIDIA A100 GPUs, which takes around 5 days.
In the video-text pre-training model, since the video representation is obtained by inputting four frames of the video into the image encoder in our framework, we have completely followed the settings of all image-text pre-training models. 
This model is also pre-trained for 100K steps on 8 NVIDIA A100 GPUs, which takes about 8 days.


\subsection{Fine-tuning settings}
\label{sec:app_finetune_setting}

For the downstream tasks, we perform grid searches over the learning rates and image resolutions. 
Finally, the learning rate of the three models is uniformly set to 1e-5, the image resolution is 384, and the max text length is 50.


\section{More Empirical Results}
\label{sec:app_additional_exp_res}

\subsection{Unimodal performance in genre classification}
\label{sec:app_unimodal_genre_cls}

\begin{table*}
\centering
\begin{tabular}{llcccccc}
\toprule
\multirow{2}{*}{\textbf{Modality}} & \multirow{2}{*}{\textbf{Source}} & \multicolumn{3}{c}{\textbf{Micro}} & \multicolumn{3}{c}{\textbf{Macro}} \\
&  &  P & R & F1 & P & R & F1 \\
\midrule
\multirow{2}{*}{Text} & BSC & 80.2 & 70.3 & 75.0 & 77.4 & 59.1 & 65.2 \\
 & Plot & 88.1 & 88.5 & 88.3 & 78.9 & 79.4 & 79.1 \\
\midrule
Image & Frame & 43.8 & 51.6 & 47.7 & 39.4 & 42.5 & 40.9 \\
\midrule
Video & Clip & 49.3 & 58.6 & 58.7 & 41.3 & 44.6 & 42.9 \\
\bottomrule
\end{tabular}
\caption{\label{tab:app_genre_cls}
The results for unimodal data in genre classification of \TVD.
}
\end{table*}

Tab.~\ref{tab:app_genre_cls} lists the results of genre classification based on TV drama data of different modalities.
Observing the performance gap between data with different modalities reveals that the modality of the data has a great impact on genre classification when processing the same plot clip.
Specifically, in the genre classification task, the performance of visual signals is significantly lower than that of text signals such as plots or BSCs.
On the one hand, this is because the content of the text signal has been manually edited, which can summarize the plot characteristics with high quality, and can more directly reflect the genre of TV dramas compared with the visual signal.
On the other hand, existing text encoders generally outperform visual encoders in semantic representation capabilities.

There are also differences in the performance of text signals from the three different sources.
Among them, TV drama abstract has the worst support for genre prediction, mainly due to the small number of tag samples in some genres.
This directly makes it difficult for the model to fully capture the genre characteristics in the case of few shots, making some categories perform poorly, resulting in the Macro results being significantly lower than the Micro results.
Genre classification based on BSC and plot text shows high performance, indicating the data value of the above text in plot analysis.
In addition, the genre classification performance based on BSC is close to that of manually edited plots in Precision and far lower in Recall, indicating that BSC content can guide the understanding of genres to a certain extent, but some noise samples lead to the inability to recall all samples.


\subsection{The number of frames in the video}
\label{sec:app_number_frame}

\begin{table*}[t]
\centering
\begin{tabular}{c|cccccc}
\toprule
\textbf{\#frames} & \textbf{VR@1} & \textbf{VR@5} & \textbf{VR@10} & \textbf{TR@1} & \textbf{TR@5} & \textbf{TR@10} \\
\midrule
 2 & 1.8 & 7.8 & 14.6 & 1.8 & 7.8 & 15.8 \\
 4 & 1.8 & 8.9 & 15.5 & 2.3 & 9.4 & 16.9 \\
 6 & 2.0 & 10.5 & 18.0 & 2.4 & 12.2 & 19.4 \\
\bottomrule
\end{tabular}
\caption{\label{tab:frames}
The video-text retrieval result of \TVD with different numbers of frames.
}
\end{table*}

In this paper, we achieve video feature acquisition by uniformly intercepting frames in the video and stitching them in the ViT.
In order to consider the quality of the video feature and its performance in downstream tasks with the different numbers of intercepted frames, we set different intercepted numbers and compared them in cross-modal retrieval. 
The results are shown in Tab.~\ref{tab:frames}.
Through observation, it can be found that with the increase in the number of intercepted frames in the video, the performance of the model has been greatly improved in terms of VR and TR indicators.
This shows that when processing video, considering as much video signal information as possible can effectively improve its performance in downstream tasks.


\subsection{Bullet screen comment generation}
\label{sec:app_gt}

\begin{table*}
\centering
\scriptsize
\begin{tabular}{clccccc}
\toprule
\textbf{\#BSC for reference} & \textbf{Settings} & \textbf{BLEU@1} & \textbf{BLEU@2} & \textbf{METEOR} & \textbf{CIDEr} \\
\midrule
\multirow{2}{*}{BSC-1} & Complete model & 9.1 & 3.9 & 4.3 & 7.6 \\
& \textit{-w/o BSC pre-train}  & 10.1 & 4.3 & 4.4 & 8.5 \\
\midrule
\multirow{2}{*}{BSC-3} & Complete model & 25.5 & 11.8 & 7.2 & 5.9 \\
& \textit{-w/o BSC pre-train}  & 24.7 & 11.0 & 7.2 & 6.5 \\
\midrule
\multirow{2}{*}{BSC-5} & Complete model & 36.6 & 17.8 & 8.8 & 5.4 \\
& \textit{-w/o BSC pre-train}  & 32.6 & 15.6 & 8.7 & 5.7 \\
\bottomrule
\end{tabular}
\caption{\label{tab:BSC_generation}
The results of BSC generation.
}
\end{table*}

\begin{table*}
\centering
\tiny
\begin{tabular}{cll}
\toprule
\textbf{Video} & \multicolumn{2}{c}{\textbf{Text}} \\
\midrule
\multirow{4}{*}{\includegraphics[width=0.3\linewidth]{figure/fig_gen_4.png}} & \multirow{2}{*}{GT:} & \begin{CJK*}{UTF8}{gbsn}塞勒斯到艾琳娜的学校找到了艾琳娜，他以斯特凡的身份和艾琳娜谈论了达蒙。\end{CJK*} \\
 & & Cyrus went to Elena's school to find Elena, and he talked to Elena about Damon as Stefan. \\
 & \multirow{2}{*}{GR:} & \begin{CJK*}{UTF8}{gbsn}斯特凡到艾琳娜那里去找艾琳娜，艾琳娜告诉他自己不会有任何事情。\end{CJK*} \\
 & & Stefan went to find Elena, and Elena told him that nothing would happen to her. \\
 & & \\
\multirow{4}{*}{\includegraphics[width=0.3\linewidth]{figure/fig_gen_5.png}} & \multirow{2}{*}{GT:} & \begin{CJK*}{UTF8}{gbsn}谢恩找到了邦妮，他要邦妮小心复活的塞拉斯。\end{CJK*}  \\
 & & Shane finds Bonnie and tells her to watch out for the resurrected Silas. \\
 & \multirow{2}{*}{GR:} & \begin{CJK*}{UTF8}{gbsn}艾琳娜和达蒙见到了马特，并要他们多注意。\end{CJK*} \\
 & & Elena and Damon meet Matt and ask them to pay more attention. \\
 & & \\
\multirow{4}{*}{\includegraphics[width=0.3\linewidth]{figure/fig_gen_6.png}} & \multirow{2}{*}{GT:} & \begin{CJK*}{UTF8}{gbsn}在那里麦克斯和凯若琳遇到了她们英俊的邻居。\end{CJK*}  \\
 & & There Max and Caroline meet their handsome neighbor. \\
 & \multirow{2}{*}{GR:} & \begin{CJK*}{UTF8}{gbsn}麦克斯和凯若琳在餐厅见面，她们正在准备约会。\end{CJK*} \\
 & & Max and Caroline meet at a restaurant. They're getting ready for a date. \\
 & & \\
\multirow{4}{*}{\includegraphics[width=0.3\linewidth]{figure/fig_gen_7.png}} & \multirow{2}{*}{GT:} & \begin{CJK*}{UTF8}{gbsn}这点让凯若琳很失望，她准备去找吃的。\end{CJK*}  \\
 & & This disappointed Caroline, she was going to find something to eat. \\
 & \multirow{2}{*}{GR:} & \begin{CJK*}{UTF8}{gbsn}凯若琳表示自己很不高兴，她要麦克斯一起出去。\end{CJK*} \\
 & & Caroline says she's upset and wants Max to go out with her. \\
\bottomrule
\end{tabular}
\caption{\label{tab:add_plot_generation_case}
More cases of plot text generation in \TVD. GT is short for ground truth. GR is short for generation results.
}
\end{table*}

In order to verify the performance of BSC pre-training in BSC generation and examine the results of introducing BSC at the same time point as evaluation references, we conducted relevant experiments, and the results are shown in Tab.~\ref{tab:BSC_generation}.
It can be found that when more BSC at the same time point are gradually introduced, the performance of the model is gradually improved, and BSC pre-training also shows a gain effect.


\subsection{More generation results}
\label{sec:app_generation_res}



\begin{figure}[ht]
\centering
\includegraphics[width=0.6\textwidth]{figure/fig_examples.pdf}
\caption{
An example of plot text generation results for multimodal LLMs.
}
\label{fig:mmllm}
\end{figure}

We list more plot text generation results here, as shown in Tab.~\ref{tab:add_plot_generation_case}.
Through observation, it can be clearly found that the video-text pre-training model, which can perform well on the public website dataset~\citep{pretrainvl12020arxiv}, does not achieve ideal results in plot generation.
This is mainly due to the difference between the TV drama data and the publicly collected dataset, i.e., the TV drama data contain more explicit entity names (first two rows in Tab.~\ref{tab:add_plot_generation_case}) and long-distance temporal relationships (last two rows in Tab.\ref{tab:add_plot_generation_case}), and the existing model is difficult to deal with such information.
Certainly, we can also find that the video-text pre-training model we used can generate more semantic information with situation context and mental state based on visual signals, indicating that the existing multimodal understanding model can try to understand the information hidden behind visual signals.

Recently, a large number of multimodal large language models (LLMs) (e.g., LLaVA\footnote{\url{https://llava.hliu.cc/}}, MiniGPT-4\footnote{\url{https://huggingface.co/spaces/Vision-CAIR/minigpt4}}, visualGLM\footnote{\url{https://huggingface.co/spaces/lykeven/visualglm-6b}}, and mPLUG-Owl\footnote{\url{https://modelscope.cn/studios/damo/mPLUG-Owl/summary}}) have achieved outstanding performance on multimodal data understanding tasks (especially images, and videos).
We selected several well-known open-source multimodal LLMs to perform the task of plot text generation in zero-shot setting. The experimental results are shown in Fig.~\ref{fig:mmllm}.
The generation results are not satisfactory. Most of the models pay more attention to the captions of the input visual signal. This may be due to the problem of understanding the task instructions, and the visual signal processing and thought chain of the existing multi-modal LLMs inferring ability is not strong enough.
Therefore, the plot comprehension and generation tasks are still challenging tasks.

% \begin{table*}
% \centering
% \begin{tabular}{cc}
% \toprule
% \textbf{Multimodal LLM} & \textbf{Example} \\
% \midrule
% LLaVA & \includegraphics[width=0.8\linewidth]{figure/llava.png} \\
% \midrule
% MiniGPT-4 & \includegraphics[width=0.8\linewidth]{figure/minigpt4.png} \\
% \midrule
% visualGLM & \includegraphics[width=0.8\linewidth]{neurips_data_2023/figure/visualglm.png} \\
% \midrule
% mPLUG-Owl & \includegraphics[width=0.8\linewidth]{neurips_data_2023/figure/mplug-owl.png} \\
% \bottomrule
% \end{tabular}
% \caption{An example of plot text generation results for multimodal LLMs.}
% \label{tab:mmllm}
% \end{table*}


\section{Limitations}
\label{sec:limitations}
\paragraph{Dataset.}
Due to the high cost of collecting and annotating data, there are currently only 83 TV dramas in \TVD.
As analysed in original paper, the genre distribution of \TVD is unbalanced.
The text data of \TVD (i.e., the plot text and Bullet Screen Comments) is in Chinese only, and heavily reflect Chinese culture and expression habits. Therefore, \TVD may exhibit gaps between real-world data from other social backgrounds.
Also, although we adhere to strict standards in our data collection pipeline, we are not able to manually check all the samples, so \TVD may contain a small percentage of annotation errors.

\paragraph{Framework.}
To reduce the environmental impact and to highlight scalability, the framework implemented in this study adopts the very popular techniques rather than the most advanced ones.
As identified in the original paper, our framework may aggregate existing distribution bias in the dataset, which should be treated very carefully.
In the task of Plot Text Generation, we did not conduct manual evaluation as our budget has run out during the creation of \TVD (annotation expense, copyright overhead, etc.).
However, we should remind readers that human judgement is still needed in some cases.

\paragraph{Experiments.}
Our framework, together with \TVD, could be used to explore a wide range of tasks, but due to time limit we only tested three most representative ones.
We made several counter-intuition observations. Although we have provided reasonable guesses on their reasons, we admit a very in-depth investigation still requires a considerable amount of work.

In general, despite the above limitations, we believe that our dataset, our framework, and our empirical findings make an important addition to the multi-modal learning, or more generally, the machine intelligence research field. 

\section{Ethics Statement}
\label{sec:ethics}

\paragraph{Data life-cycle.}
Our data annotation and release plan has been examined and approved by relevant institutional committees.
All annotators participating the creation of \TVD have agreed relevant terms and received a set of pre-annotation training.
They were provided with remuneration that is significantly higher than the local minimum wage.
At different stages of data collection, we complied with corresponding legal and ethical requirements.
% Specifically, our data collection belongs to the Personal Information Protection Law of the People's Republic of China.
In order to protect the privacy of TV drama viewer, we have also anonymised all data (especially Bullet Screen Comments), ensuring no information that can uniquely identify users is excluded from our dataset and research process.
We have obtained all data copyrights relevant to this paper.
We have prepared detailed instructions, licenses, and data use agreement for future users of \TVD, aiming to minimise the possibility of any kind of misuse (e.g., to train a malicious Bullet Screen Comment writer).



\paragraph{Data bias.}
Most of the TV dramas included in \TVD are produced in a fictional context, and the relevant content may have the specific bias introduced by the original creators.
Meanwhile, the filming techniques and editing strategies commonly used in film and television photography will also highlight some dramatic plots and may have a preference for some content.
In \TVD,  the Bullet Screen Comments carry the opinions of the user, which may therefore contain personal bias.
The above bias may be inherited or even aggregated by the downstream algorithms.
Thus, we urge future users of \TVD to pay special  attention to this matter.

